\section{The Second Fundamental Form}

Let $(\widetilde M,\widetilde g)$ be an $m$-dimensional Riemannian or
pseudo-Riemannian manifold without boundary, and let $(M,g)$ be an embedded
$n$-dimensional Riemannian submanifold without boundary, with
$g=\iota_M^*\widetilde g$. All statements are local and therefore also apply
to isometric immersions after restriction to embedded neighborhoods. We
restrict to positive-definite induced metrics; mixed-signature submanifolds
require additional qualifications (see \cite{dC92}).

Ambient covariant derivatives and curvatures carry tildes; untilded symbols
refer to the induced geometry of $M$. The orthogonal splitting
$T\widetilde M|_M=TM\oplus NM$ has smooth projections
\[
\pi^\top:T\widetilde M|_M\to TM,
\qquad
\pi^\perp:T\widetilde M|_M\to NM.
\]
For a section $A$ of $T\widetilde M|_M$, write
$A^\top=\pi^\top A$ and $A^\perp=\pi^\perp A$. If
$X,Y\in\mathfrak X(M)$ are extended to ambient vector fields, then
\begin{equation}
\widetilde\nabla_XY
=(\widetilde\nabla_XY)^\top+(\widetilde\nabla_XY)^\perp.
\tag{8.1}
\end{equation}
The \emph{second fundamental form} is the normal-bundle-valued map
\[
\mathrm{II}(X,Y)=(\widetilde\nabla_XY)^\perp.
\]

\begin{proposition}[Properties of the Second Fundamental Form]
\label{prop:second-fundamental-form-properties}
\dcref{ch8:8.1}

Let $(M,g)$ be an embedded Riemannian submanifold of a Riemannian or
pseudo-Riemannian manifold $(\widetilde M,\widetilde g)$. For
$X,Y\in\mathfrak X(M)$:
\begin{enumerate}[label=(\alph*)]
\item $\mathrm{II}(X,Y)$ is independent of the chosen ambient extensions.
\item $\mathrm{II}$ is $C^\infty(M)$-bilinear.
\item $\mathrm{II}(X,Y)=\mathrm{II}(Y,X)$.
\item $\mathrm{II}(X,Y)|_p$ depends only on $X_p$ and $Y_p$.
\end{enumerate}
\end{proposition}

\begin{proof}
\sketch For fixed ambient extensions, symmetry of $\widetilde\nabla$ gives
\[
\mathrm{II}(X,Y)-\mathrm{II}(Y,X)
=(\widetilde\nabla_XY-\widetilde\nabla_YX)^\perp=[X,Y]^\perp=0,
\]
because the bracket of fields tangent along $M$ is tangent to $M$. At
$p\in M$, $\widetilde\nabla_XY|_p$ depends only on $X_p$, proving independence
of the extension of $X$ and $C^\infty(M)$-linearity in $X$. Symmetry gives the
same conclusions for $Y$, and hence pointwise dependence on both arguments.
\end{proof}

Thus $\mathrm{II}(v,w)$ is well defined for $v,w\in T_pM$ by extending them
to tangent vector fields.

\begin{theorem}[The Gauss Formula]
\label{thm:gauss-formula}
\dcref{ch8:8.2}
\uses{prop:second-fundamental-form-properties,ex:tangential-connection-submanifold,prop:tangential-connection-compatible-metric,prop:tangential-connection-symmetric}

Under the preceding hypotheses, for arbitrary ambient extensions of
$X,Y\in\mathfrak X(M)$,
\[
\widetilde\nabla_XY=\nabla_XY+\mathrm{II}(X,Y)
\qquad\text{along }M.
\]
\end{theorem}

\begin{proof}
\sketch Define
$\nabla_X^\top Y=(\widetilde\nabla_XY)^\top$. The extension-independence
argument for the second fundamental form shows that $\nabla^\top$ is a
connection on $M$. Ambient metric compatibility and orthogonal projection
make it compatible with $g$, while ambient symmetry and tangency of brackets
make it symmetric. Uniqueness of the Levi-Civita connection gives
$\nabla^\top=\nabla$. Equation (8.1) then proves the formula.
\end{proof}

\begin{corollary}[The Gauss Formula Along a Curve]
\label{cor:gauss-formula-along-curve}
\dcref{ch8:8.3}
\uses{thm:gauss-formula}

Let $\gamma:I\to M$ be smooth and let $X$ be a smooth vector field along
$\gamma$ tangent to $M$. Its ambient and intrinsic covariant derivatives
satisfy
\begin{equation}
\widetilde D_tX=D_tX+\mathrm{II}(\gamma',X).
\tag{8.2}
\end{equation}
\end{corollary}

\begin{proof}
\sketch Near $\gamma(t_0)$ choose an adapted orthonormal frame
$(E_1,\ldots,E_m)$ and write $X=X^iE_i$ for $1\leq i\leq n$. The product rule
and Gauss formula give
\[
\widetilde D_tX
=\sum_{i=1}^n\left(\dot X^iE_i+X^i\nabla_{\gamma'}E_i
+X^i\mathrm{II}(\gamma',E_i)\right)
=D_tX+\mathrm{II}(\gamma',X).
\]
\end{proof}

For $N\in\Gamma(NM)$, define the symmetric scalar-valued form
\begin{equation}
\mathrm{II}_N(X,Y)=\langle N,\mathrm{II}(X,Y)\rangle.
\tag{8.3}
\end{equation}
Its associated self-adjoint bundle map $W_N:TM\to TM$, the
\emph{Weingarten map in the direction $N$}, is characterized by
\begin{equation}
\langle W_N(X),Y\rangle
=\mathrm{II}_N(X,Y)
=\langle N,\mathrm{II}(X,Y)\rangle.
\tag{8.4}
\end{equation}

\begin{proposition}[The Weingarten Equation]
\label{prop:weingarten-equation}
\dcref{ch8:8.4}

For every $X\in\mathfrak X(M)$ and $N\in\Gamma(NM)$,
\begin{equation}
(\widetilde\nabla_XN)^\top=-W_N(X),
\tag{8.5}
\end{equation}
where $N$ is extended arbitrarily to a neighborhood of $M$ in
$\widetilde M$.
\end{proposition}

\begin{proof}
\sketch For arbitrary $Y\in\mathfrak X(M)$, metric compatibility, normality
of $N$, and the Gauss formula yield
\begin{align*}
0=X\langle N,Y\rangle
&=\langle\widetilde\nabla_XN,Y\rangle
 +\langle N,\widetilde\nabla_XY\rangle\\
&=\langle\widetilde\nabla_XN,Y\rangle
 +\langle N,\mathrm{II}(X,Y)\rangle\\
&=\langle(\widetilde\nabla_XN)^\top+W_N(X),Y\rangle.
\end{align*}
Since $Y$ is arbitrary, (8.5) follows.
\end{proof}

\begin{theorem}[The Gauss Equation]
\label{thm:gauss-equation}
\dcref{ch8:8.5}
\uses{prop:covariant-derivative-depends-on-curve}

For all $W,X,Y,Z\in\mathfrak X(M)$,
\[
\widetilde{Rm}(W,X,Y,Z)
=Rm(W,X,Y,Z)
-\langle\mathrm{II}(W,Z),\mathrm{II}(X,Y)\rangle
+\langle\mathrm{II}(W,Y),\mathrm{II}(X,Z)\rangle.
\]
\end{theorem}

\begin{proof}
\sketch Extend $W,X,Y,Z$ ambiently. Apply the Gauss formula inside the
definition of ambient curvature:
\begin{align*}
\widetilde{Rm}(W,X,Y,Z)
&=\big\langle
\widetilde\nabla_W(\nabla_XY+\mathrm{II}(X,Y))
-\widetilde\nabla_X(\nabla_WY+\mathrm{II}(W,Y))\\
&\qquad-\widetilde\nabla_{[W,X]}Y,Z
\big\rangle.
\end{align*}
The Weingarten equation converts the tangential components of the two normal
derivatives into
$-\langle\mathrm{II}(X,Y),\mathrm{II}(W,Z)\rangle$ and
$+\langle\mathrm{II}(W,Y),\mathrm{II}(X,Z)\rangle$. Applying the Gauss
formula to the remaining tangential derivatives leaves
$\langle R(W,X)Y,Z\rangle$, proving the equation.
\end{proof}

The \emph{normal connection} is
\[
\nabla_X^\perp N=(\widetilde\nabla_XN)^\perp
\qquad
(X\in\mathfrak X(M),\ N\in\Gamma(NM)).
\]

\begin{proposition}
\label{prop:normal-connection-well-defined}
\dcref{ch8:8.6}

The operator $\nabla^\perp$ is a well-defined connection on $NM$ compatible
with the restricted ambient metric:
\[
X\langle N_1,N_2\rangle
=\langle\nabla_X^\perp N_1,N_2\rangle
 +\langle N_1,\nabla_X^\perp N_2\rangle
\]
for $N_1,N_2\in\Gamma(NM)$ and $X\in\mathfrak X(M)$.
\end{proposition}

\begin{exercise}
\label{exer:prove-normal-connection}
\dcref{ch8:8.7}
\uses{prop:normal-connection-well-defined}

Prove Proposition~\ref{prop:normal-connection-well-defined}.
\end{exercise}

Let $F\to M$ be the bundle with fiber
\[
F_p=\operatorname{Bil}(T_pM\times T_pM,N_pM).
\]
For $B\in\Gamma(F)$, define
\[
(\nabla_X^FB)(Y,Z)
=\nabla_X^\perp(B(Y,Z))
-B(\nabla_XY,Z)-B(Y,\nabla_XZ).
\]

\begin{exercise}
\label{exer:connection-in-bundle}
\dcref{ch8:8.8}

Prove that $\nabla^F$ is a connection on $F$.
\end{exercise}

\begin{theorem}[The Codazzi Equation]
\label{thm:codazzi-equation}
\dcref{ch8:8.9}
\uses{thm:gauss-equation}

For all $W,X,Y\in\mathfrak X(M)$,
\begin{equation}
(\widetilde R(W,X)Y)^\perp
=(\nabla_W^F\mathrm{II})(X,Y)-(\nabla_X^F\mathrm{II})(W,Y).
\tag{8.6}
\end{equation}
\end{theorem}

\begin{proof}
\sketch It suffices to pair both sides with an arbitrary normal field $N$:
\begin{equation}
\langle\widetilde R(W,X)Y,N\rangle
=\langle(\nabla_W^F\mathrm{II})(X,Y),N\rangle
-\langle(\nabla_X^F\mathrm{II})(W,Y),N\rangle.
\tag{8.7}
\end{equation}
After applying the Gauss formula to the ambient curvature definition, retain
only normal components. The first ambient derivative contributes
\[
\mathrm{II}(W,\nabla_XY)
+(\nabla_W^F\mathrm{II})(X,Y)
+\mathrm{II}(\nabla_WX,Y)
+\mathrm{II}(X,\nabla_WY),
\]
and the second contributes the analogous expression with $W,X$ interchanged,
while the bracket term contributes $-\mathrm{II}([W,X],Y)$. The terms
containing $\nabla_XY$ and $\nabla_WY$ cancel, and
\[
\mathrm{II}(\nabla_WX-\nabla_XW-[W,X],Y)=0
\]
by symmetry of $\nabla$. Pairing the remaining terms with $N$ gives (8.7).
\end{proof}

\section{Curvature of a Curve}

For a smooth unit-speed curve $\gamma:I\to M$ in a Riemannian or
pseudo-Riemannian manifold, its \emph{geodesic curvature} is
\[
\kappa(t)=|D_t\gamma'(t)|.
\]
For a regular curve in a Riemannian manifold, define curvature using any
unit-speed reparametrization. In the pseudo-Riemannian case, this definition
is restricted to curves satisfying $|\gamma'(t)|\neq0$ everywhere. A
unit-speed curve has zero geodesic curvature exactly when it is a geodesic.

If $M$ is a Riemannian submanifold of
$(\widetilde M,\widetilde g)$, a regular curve in $M$ has intrinsic curvature
$\kappa$ computed in $M$ and extrinsic curvature $\widetilde\kappa$ computed
in $\widetilde M$.

\begin{proposition}[Geometric Interpretation of II]
\label{prop:geometric-interpretation-ii}
\dcref{ch8:8.10}
\uses{cor:gauss-formula-along-curve}

Let $p\in M$ and $v\in T_pM$.
\begin{enumerate}
\item[(a)] $\mathrm{II}(v,v)$ is the ambient acceleration at $p$ of the
$g$-geodesic $\gamma_v$.
\item[(b)] If $v$ is a unit vector, then $|\mathrm{II}(v,v)|$ is the ambient
curvature of $\gamma_v$ at $p$.
\end{enumerate}
\end{proposition}

\begin{proof}
\sketch For any regular $\gamma:(-\varepsilon,\varepsilon)\to M$ with
$\gamma(0)=p$ and $\gamma'(0)=v$, (8.2) gives
\[
\widetilde D_t\gamma'
=D_t\gamma'+\mathrm{II}(\gamma',\gamma').
\]
For $\gamma=\gamma_v$, the intrinsic derivative vanishes. If $v$ is unit,
the ambient acceleration norm is the ambient curvature.
\end{proof}

\begin{lemma}
\label{lem:symmetric-bilinear-form-unit-vectors}
\dcref{ch8:8.11}

Let $V$ be an inner product space, $W$ a vector space, and
$B,B':V\times V\to W$ symmetric bilinear maps. If
$B(v,v)=B'(v,v)$ for every unit vector $v$, then $B=B'$.
\end{lemma}

\begin{proof}
\sketch Homogeneity extends the diagonal equality to every $v\in V$.
Polarization then gives
\[
B(v,w)=\frac14\bigl(B(v+w,v+w)-B(v-w,v-w)\bigr),
\]
and the same identity for $B'$ proves equality.
\end{proof}

A Riemannian submanifold $(M,g)$ of $(\widetilde M,\widetilde g)$ is
\emph{totally geodesic} if every ambient geodesic tangent to $M$ at a time
$t_0$ remains in $M$ on some interval about $t_0$.

\begin{proposition}
\label{prop:totally-geodesic-equivalence}
\dcref{ch8:8.12}
\uses{thm:gauss-formula,lem:symmetric-bilinear-form-unit-vectors}

The following are equivalent:
\begin{enumerate}
\item[(a)] $M$ is totally geodesic in $\widetilde M$.
\item[(b)] Every $g$-geodesic in $M$ is an ambient
$\widetilde g$-geodesic.
\item[(c)] $\mathrm{II}=0$ identically.
\end{enumerate}
\end{proposition}

\begin{proof}
\sketch Assume (a), let $\gamma$ be a $g$-geodesic, and at each $t_0$ let
$\widetilde\gamma$ be the ambient geodesic with the same initial data.
Locally $\widetilde\gamma$ lies in $M$, and (8.2) decomposes
$0=\widetilde D_t\widetilde\gamma'$ into tangent and normal parts. Thus it is
also a $g$-geodesic, so uniqueness gives
$\widetilde\gamma=\gamma$. Hence (b).

Under (b), for the $g$-geodesic with initial vector $v\in T_pM$, both
intrinsic and ambient accelerations vanish. Formula (8.2) gives
$\mathrm{II}(v,v)=0$, and
Lemma~\ref{lem:symmetric-bilinear-form-unit-vectors} gives (c).

If (c) holds, every $g$-geodesic is ambient by (8.2). Given an ambient
geodesic initially tangent to $M$, compare it with the $g$-geodesic having
the same initial data. Ambient geodesic uniqueness makes them agree near the
initial time, proving (a).
\end{proof}
\section{Hypersurfaces}

Throughout these sections, $(M,g)$ is an embedded Riemannian hypersurface of
dimension $n$ in a Riemannian manifold $(\widetilde M,\widetilde g)$ of
dimension $n+1$.

\section{The Scalar Second Fundamental Form and the Shape Operator}

Locally, choose a smooth unit normal field $N$ along $M$. Such a choice always
exists near each point; if both $M$ and $\widetilde M$ are orientable, their
orientations determine a global choice. The \emph{scalar second fundamental
form} relative to $N$ is the symmetric covariant tensor
\begin{equation}
h(X,Y)=\langle N,\operatorname{II}(X,Y)\rangle. \tag{8.8}
\end{equation}
The Gauss formula and orthogonality give
\begin{equation}
h(X,Y)=\langle N,\widetilde\nabla_XY\rangle, \tag{8.9}
\end{equation}
and, because $N$ spans the normal bundle,
\begin{equation}
\operatorname{II}(X,Y)=h(X,Y)N. \tag{8.10}
\end{equation}
Replacing $N$ by $-N$ changes $h$ to $-h$.

The \emph{shape operator} $s=W_N$ is obtained from $h$ by raising one index:
\[
\langle sX,Y\rangle=h(X,Y).
\]
It is self-adjoint, and replacing $N$ by $-N$ changes $s$ to $-s$.

For symmetric covariant $2$-tensors $h,k$, use the Kulkarni--Nomizu product
\begin{align*}
(h\owedge k)(w,x,y,z)
&=h(w,z)k(x,y)+h(x,y)k(w,z)\\
&\quad-h(w,y)k(x,z)-h(x,z)k(w,y),
\end{align*}
and for a symmetric covariant $2$-tensor $T$, use the exterior covariant
derivative
\[
(DT)(x,y,z)=-(\nabla T)(x,y,z)+(\nabla T)(x,z,y).
\]

\begin{theorem}[Fundamental Equations for a Hypersurface]
\label{thm:fundamental-equations-hypersurface}
\dcref{ch8:8.13}
Let $(M,g)$ be a Riemannian hypersurface in $(\overline M,\overline g)$, and
let $N$ be a smooth unit normal field along $M$.
\begin{enumerate}
\item[(a)] If $X,Y\in\mathfrak X(M)$ are extended to an open subset of
$\overline M$, then
\[
\overline\nabla_XY=\nabla_XY+h(X,Y)N.
\]
\item[(b)] If $\gamma\colon I\to M$ is smooth and $X$ is a smooth tangent
vector field along $\gamma$, then
\[
\overline D_tX=D_tX+h(\gamma',X)N.
\]
\item[(c)] For every $X\in\mathfrak X(M)$,
\begin{equation}
\overline\nabla_XN=-sX. \tag{8.11}
\end{equation}
\item[(d)] For all $W,X,Y,Z\in\mathfrak X(M)$,
\[
\overline{\operatorname{Rm}}(W,X,Y,Z)
=\operatorname{Rm}(W,X,Y,Z)-\frac12(h\owedge h)(W,X,Y,Z).
\]
\item[(e)] For all $W,X,Y\in\mathfrak X(M)$,
\begin{equation}
\overline{\operatorname{Rm}}(W,X,Y,N)=(Dh)(Y,W,X). \tag{8.12}
\end{equation}
\end{enumerate}
\end{theorem}

\begin{proof}
\sketch Substitution of (8.10) into the general Gauss formulas and Gauss
equation gives (a), (b), and (d). The general Weingarten equation gives
$(\overline\nabla_XN)^\top=-sX$, while
\[
\langle\overline\nabla_XN,N\rangle
=\frac12X(|N|^2)=0;
\]
hence $\overline\nabla_XN$ is tangent and (c) follows.

For (e), the unit-length condition gives
$\nabla_X^\perp N=0$, since $N$ spans the normal bundle. Therefore
\begin{align*}
(\nabla_W^\perp\operatorname{II})(X,Y)
&=\nabla_W^\perp(h(X,Y)N)
  -\operatorname{II}(\nabla_WX,Y)
  -\operatorname{II}(X,\nabla_WY)\\
&=(\nabla_Wh)(X,Y)N.
\end{align*}
The general Codazzi equation then yields
\begin{align*}
\overline{\operatorname{Rm}}(W,X,Y,N)
&=(\nabla_Wh)(X,Y)-(\nabla_Xh)(W,Y)\\
&=(\nabla h)(Y,X,W)-(\nabla h)(Y,W,X)\\
&=(Dh)(Y,W,X),
\end{align*}
using symmetry of $\nabla h$ in its first two tensor arguments.
\end{proof}

\section{Principal Curvatures}

\begin{lemma}
\label{lem:self-adjoint-eigenvector-extremum}
\dcref{ch8:8.14}
Let $V$ be a finite-dimensional inner product space, let
$s\colon V\to V$ be self-adjoint, and let $C$ be the unit sphere in $V$.
The function $v\mapsto\langle sv,v\rangle$ attains its maximum at some
$v_0\in C$, and every maximizer is an eigenvector of $s$ with eigenvalue
$\lambda_0=\langle sv_0,v_0\rangle$.
\end{lemma}

\begin{exercise}
\label{exer:lagrange-multiplier-eigenvector}
\dcref{ch8:8.15}
\uses{prop:lagrange-multiplier-rule}
Prove the preceding lemma using the Lagrange multiplier rule.
\end{exercise}

\begin{proposition}[Finite-Dimensional Spectral Theorem]
\label{prop:finite-dimensional-spectral-theorem}
\dcref{ch8:8.16}
\uses{lem:self-adjoint-eigenvector-extremum}
Every self-adjoint endomorphism of a finite-dimensional inner product space
has an orthonormal basis of eigenvectors, and all its eigenvalues are real.
\end{proposition}

\begin{proof}
\sketch Induct on $n=\dim V$. The result is immediate for $n=1$. In dimension
$n+1$, the preceding lemma supplies a unit eigenvector $b_0$ with real
eigenvalue. For $B=\operatorname{span}\{b_0\}$, self-adjointness implies
$s(B^\perp)\subseteq B^\perp$. Apply the induction hypothesis to
$s|_{B^\perp}$ and adjoin $b_0$ to the resulting orthonormal eigenbasis.
\end{proof}

At $p\in M$, let $\kappa_1,\ldots,\kappa_n$ be the eigenvalues of the shape
operator. In an orthonormal eigenbasis $(b_i)$,
\[
sb_i=\kappa_i b_i,
\qquad
h(v,w)=\sum_{i=1}^n\kappa_i v^iw^i.
\]
The $\kappa_i$ are the \emph{principal curvatures}, and the corresponding
eigenspaces are the \emph{principal directions}. Reversing $N$ negates every
principal curvature but leaves the principal directions unchanged.

The \emph{Gaussian curvature} and \emph{mean curvature} are
\[
K=\det s=\prod_{i=1}^n\kappa_i,
\qquad
H=\frac1n\operatorname{tr}s
=\frac1n\operatorname{tr}_g h
=\frac1n\sum_{i=1}^n\kappa_i.
\]
Under $N\mapsto-N$, the mean curvature changes sign and the Gaussian curvature
is multiplied by $(-1)^n$.

\section{Computations in Semigeodesic Coordinates}

Let $(x^1,\ldots,x^n,v)$ be semigeodesic coordinates on
$U\subseteq\widetilde M$. For each $a$ with
$M_a=v^{-1}(a)\ne\varnothing$, the level set $M_a$ is a properly embedded
hypersurface. Its induced metric is denoted by $g_a$, and
\begin{equation}
\widetilde g=dv^2+sum_{\alpha,\beta=1}^n
g_{\alpha\beta}(x^1,\ldots,x^n,v)\,dx^\alpha dx^\beta. \tag{8.13}
\end{equation}
The restrictions of $(x^1,\ldots,x^n)$ are coordinates on $M_a$,
$g_a=g_{\alpha\beta}dx^\alpha dx^\beta|_{v=a}$, and $\partial_v$ is a unit
normal along $M_a$.

\begin{proposition}
\label{prop:semigeodesic-coordinates-second-form}
\dcref{ch8:8.17}
With respect to the normal $\partial_v$, the scalar second fundamental form,
shape operator, and mean curvature of $M_a$ have components
\[
(h_a)_{\alpha\beta}
=-\frac12\partial_vg_{\alpha\beta}|_{v=a},
\qquad
(s_a)_\beta{}^\alpha
=-\frac12g^{\alpha\gamma}\partial_vg_{\gamma\beta}|_{v=a},
\qquad
H_a=-\frac1{2n}g^{\alpha\beta}\partial_vg_{\alpha\beta}|_{v=a}.
\]
\end{proposition}

\begin{proof}
\sketch In semigeodesic coordinates, the normal component of
$\widetilde\nabla_{\partial_\alpha}\partial_\beta$ is
\[
\widetilde\Gamma^{n+1}_{\alpha\beta}\partial_v
=-\frac12\partial_vg_{\alpha\beta}\,\partial_v.
\]
Equation (8.9) gives the formula for $h_a$ after restriction to $v=a$.
Raising an index with $(g^{\alpha\gamma})$ gives $s_a$, and taking
$n^{-1}$ times its trace gives $H_a$.
\end{proof}

\section{Minimal Hypersurfaces}

For a compact embedded hypersurface $M$ with nonempty boundary, \emph{area}
means its induced $n$-dimensional volume. It is \emph{area-minimizing} if its
area is least among compact embedded hypersurfaces with the same boundary.

\begin{theorem}
\label{thm:area-minimizing-zero-mean-curvature}
\dcref{ch8:8.18}
\uses{ex:fermi-coordinates-hypersurface}
Let $M$ be a compact codimension-one submanifold with nonempty boundary in a
Riemannian $(n+1)$-manifold $(\widetilde M,\widetilde g)$. If $M$ is
area-minimizing, then its mean curvature vanishes identically.
\end{theorem}

\begin{proof}
\sketch Fix $p\in\operatorname{Int}M$ and choose Fermi coordinates
$(x^1,\ldots,x^n,v)$ on $U\subseteq\widetilde M$ such that
$\widetilde U=U\cap M$ is a coordinate ball disjoint from $\partial M$.
For arbitrary $\varphi\in C_c^\infty(\widetilde U)$ and sufficiently small
$t$, replace $M\cap U$ by the graph $v=t\varphi(x)$ to obtain $M_t$. The graph
map
\[
f_t(x)=(x^1,\ldots,x^n,t\varphi(x))
\]
extends by the identity to a diffeomorphism $F_t\colon M\to M_t$. Pull the
induced metric back to $M$ and denote it by $g_t$. On $\widetilde U$,
\begin{equation}
(g_t)_{\alpha\beta}
=t^2\partial_\alpha\varphi\,\partial_\beta\varphi
+g_{\alpha\beta}(x,t\varphi(x)), \tag{8.15}
\end{equation}
and $g_t=g$ outside $\widetilde U$.

Since the area has a minimum at $t=0$, compute its first variation:
\begin{equation}
\left.\frac d{dt}\right|_{t=0}\operatorname{Area}(M_t)
=\int_{\widetilde U}\frac12(\det g)^{-1/2}
 \left.\frac\partial{\partial t}\right|_{t=0}\det g_t
 \,dx^1\cdots dx^n. \tag{8.16}
\end{equation}
Writing $\operatorname{cof}^{\alpha\beta}$ for the matrix cofactors gives
\begin{equation}
\left.\frac\partial{\partial t}\right|_{t=0}\det g_t
=\operatorname{cof}^{\alpha\beta}
 \left.\frac\partial{\partial t}\right|_{t=0}(g_t)_{\alpha\beta}. \tag{8.17}
\end{equation}
By Cramer's rule and (8.15), this equals
\[
(\det g)g^{\alpha\beta}\partial_vg_{\alpha\beta}\,\varphi.
\]
The semigeodesic formula for $H$ therefore yields
\begin{equation}
\left.\frac d{dt}\right|_{t=0}\operatorname{Area}(M_t)
=-n\int_{\widetilde U}H\varphi\,dV_g. \tag{8.18}
\end{equation}
The left side is zero for every compactly supported $\varphi$. A nonnegative
bump function supported where $H$ has either fixed positive or fixed negative
sign rules out $H(p)\ne0$. Thus $H=0$ on the interior, and continuity gives
$H=0$ on all of $M$.
\end{proof}

An immersed or embedded hypersurface, with or without boundary, is
\emph{minimal} if $H\equiv0$. This condition means stationarity for area, not
global area minimization, although sufficiently small pieces of a minimal
hypersurface are area-minimizing.

\begin{theorem}
\label{thm:constant-mean-curvature-isoperimetric}
\dcref{ch8:8.19}
\uses{thm:area-minimizing-zero-mean-curvature,thm:tubular-neighborhood}
Let $(\widetilde M,\widetilde g)$ be a Riemannian $(n+1)$-manifold, let
$D\subseteq\widetilde M$ be a compact regular domain, and set $M=\partial D$.
If $M$ minimizes surface area among boundaries of compact regular domains with
the same volume as $D$, then its mean curvature with respect to the outward
unit normal is constant.
\end{theorem}

\begin{proof}
\sketch Suppose $H$ is not constant and choose $p,q\in M$ with
$H(p)<H(q)$. By compactness and the tubular-neighborhood theorem, choose
disjoint ambient Fermi coordinate neighborhoods $\widetilde U$ of $p$ and
$\widetilde W$ of $q$, and set
$U=M\cap\widetilde U$ and $W=M\cap\widetilde W$. Choose normal coordinates
$v,w$ increasing outward, so locally $D$ is given by $v\leq0$ and $w\leq0$.
Shrink the neighborhoods and choose constants
\[
H(p)<H_1<H_2<H(q),
\qquad H\leq H_1\text{ on }U,
\qquad H\geq H_2\text{ on }W.
\]
Choose nonnegative functions
$\varphi\in C_c^\infty(U)$ and $\psi\in C_c^\infty(W)$ normalized by
\[
\int_U\varphi\,dV_g=\int_W\psi\,dV_g=1.
\]

For small $s,t$, perturb the two boundary patches to the graphs
$v=s\varphi$ and $w=t\psi$ by setting
\begin{align*}
D_{s,t}
&=\{z\in\widetilde U:v(z)\leq s\varphi(x(z))\}\\
&\quad\cup\{z\in\widetilde W:w(z)\leq t\psi(y(z))\}\\
&\quad\cup\bigl(D\setminus(\widetilde U\cup\widetilde W)\bigr),
\qquad M_{s,t}=\partial D_{s,t}.
\end{align*}
Put
\[
V(s,t)=\operatorname{Vol}(D_{s,t}),
\qquad A(s,t)=\operatorname{Area}(M_{s,t}).
\]
The first-variation calculation (8.18) gives
\[
\frac{\partial A}{\partial s}(0,0)
=-n\int_UH\varphi\,dV_g,
\qquad
\frac{\partial A}{\partial t}(0,0)
=-n\int_WH\psi\,dV_g.
\]
The fundamental theorem of calculus in the normal coordinates gives
\[
\frac{\partial V}{\partial s}(0,0)
=\int_U\varphi\,dV_g=1,
\qquad
\frac{\partial V}{\partial t}(0,0)
=\int_W\psi\,dV_g=1.
\]
The implicit function theorem supplies a smooth $\lambda(s)$ with
$V(s,\lambda(s))=V(0,0)$ and $\lambda'(0)=-1$. Area minimality under this
volume-preserving variation implies
\[
0=\left.\frac d{ds}\right|_{s=0}A(s,\lambda(s))
=-n\int_UH\varphi\,dV_g+n\int_WH\psi\,dV_g.
\]
But nonnegativity and normalization of the test functions give
\[
\int_UH\varphi\,dV_g
\leq H_1<H_2
\leq\int_WH\psi\,dV_g,
\]
a contradiction.
\end{proof}

For further treatment of minimal and constant-mean-curvature hypersurfaces,
see \cite{CM11}.
\section{Hypersurfaces in Euclidean Space}

Let $M\subseteq\mathbb R^{n+1}$ be an embedded hypersurface with induced
metric $g$, and let $h$ be its scalar second fundamental form for a chosen unit
normal.  Since the ambient Euclidean curvature vanishes, the Gauss and Codazzi
equations are
\begin{equation}
\tfrac12h\owedge h=\mathrm{Rm}, \tag{8.19}
\end{equation}
\begin{equation}
Dh=0, \tag{8.20}
\end{equation}
or, in a local frame,
\begin{equation}
h_{il}h_{jk}-h_{ik}h_{jl}=R_{ijkl}, \tag{8.21}
\end{equation}
\begin{equation}
h_{ij;k}-h_{ik;j}=0. \tag{8.22}
\end{equation}
A symmetric $2$-tensor satisfying $Dh=0$ is a \emph{Codazzi tensor}.

\begin{exercise}
\label{exer:codazzi-tensor-characterization}
\dcref{ch8:8.20}
Show that a smooth $2$-tensor $h$ is a Codazzi tensor if and only if both $h$
and $\nabla h$ are symmetric.
\end{exercise}

Conversely, (8.19)--(8.20) are locally sufficient: if $g$ is a Riemannian
metric and $h$ is symmetric and satisfies these equations, then every point
has a neighborhood admitting an isometric immersion into $\mathbb R^{n+1}$
whose scalar second fundamental form is $h$.

For $v\in T_pM$ with $|v|=1$, let $\gamma_v$ be the intrinsic geodesic with
initial velocity $v$.  Its Euclidean acceleration is
\[
\gamma_v''(0)=h(v,v)N_p.
\]
Thus $h(v,v)=\langle\gamma_v''(0),N_p\rangle$, with sign determined by the
chosen normal.

\begin{proposition}[Osculating Circle]
\label{prop:osculating-circle}
\dcref{ch8:8.21}
Let $\gamma:I\to\mathbb R^m$ be unit speed, let $t_0\in I$, and assume
$\kappa(t_0)\ne0$.
\begin{enumerate}
\item[(a)] There is a unique unit-speed circle
$c:\mathbb R\to\mathbb R^m$ whose position, velocity, and acceleration at
$t_0$ agree with those of $\gamma$.  It is the \emph{osculating circle} at
$\gamma(t_0)$.
\item[(b)] If its radius is $R$, then $\kappa(t_0)=1/R$.
\end{enumerate}
\end{proposition}

\begin{proof}
\sketch
For orthonormal $v,w\in\mathbb R^m$, the circle with center $q$ and radius
$R$ has unit-speed parametrization
\[
c(t)=q+R\cos\left(\frac{t-t_0}{R}\right)v
+R\sin\left(\frac{t-t_0}{R}\right)w,
\]
and
\[
c(t_0)=q+Rv,
\qquad c'(t_0)=w,
\qquad c''(t_0)=-\frac1Rv.
\]
Taking
\[
R=\frac1{|\gamma''(t_0)|}=\frac1{\kappa(t_0)},
\quad v=-R\gamma''(t_0),
\quad w=\gamma'(t_0),
\quad q=\gamma(t_0)-Rv
\]
gives the required circle and radius.
\end{proof}

\begin{exercise}
\label{exer:osculating-circle-uniqueness}
\dcref{ch8:8.22}
Complete the preceding proof by establishing uniqueness of the osculating
circle.
\end{exercise}

\section{Computations in Euclidean Space}

Let $X:U\to M\subseteq\mathbb R^{n+1}$ be a local parametrization, with
$X_i=\partial_iX$.  The vectors $X_1,\ldots,X_n$ span $TM$, and Gram--Schmidt
against any independent ambient coordinate vector produces either choice of
unit normal $N$.

\begin{proposition}[Coordinate Formula for the Second Fundamental Form]
\label{prop:second-fundamental-form-coordinate-formula}
\dcref{ch8:8.23}
For a unit normal field $N$ along the parametrized hypersurface,
\begin{equation}
h(X_i,X_j)=\left\langle
\frac{\partial^2X}{\partial u^i\partial u^j},N\right\rangle. \tag{8.23}
\end{equation}
\end{proposition}

\begin{proof}
\sketch
The identity $\langle X_j,N\rangle=0$ implies
\[
0=\left\langle\partial_i\partial_jX,N\right\rangle
+\left\langle X_j,\overline\nabla_{X_i}N\right\rangle.
\]
The Weingarten equation gives
$\overline\nabla_{X_i}N=-s(X_i)$, so the second term is
$-h(X_j,X_i)$, proving (8.23).
\end{proof}

\begin{proposition}[Principal Curvatures and Local Shape]
\label{prop:principal-curvatures-local-shape}
\dcref{ch8:8.24}
Let $M\subset\mathbb R^{n+1}$ be a hypersurface, let $p\in M$, and let
$k_1,\ldots,k_n$ be its principal curvatures at $p$ for a chosen unit normal.
Some Euclidean isometry sends $p$ to $0$ and a neighborhood of $p$ in $M$ to
a graph $x^{n+1}=f(x^1,\ldots,x^n)$ with
\begin{equation}
f(x)=\frac12\left(k_1(x^1)^2+\cdots+k_n(x^n)^2\right)+O(|x|^3). \tag{8.24}
\end{equation}
\end{proposition}

\begin{proof}
\sketch
Translate and rotate so that $p=0$, $T_pM$ is the first-coordinate
hyperplane, the selected normal is $\partial_{n+1}$, and the second
fundamental form is diagonal with entries $k_1,\ldots,k_n$.  The implicit
function theorem represents $M$ locally by
$X(u)=(u^1,\ldots,u^n,f(u))$, with $f(0)=0$ and $df(0)=0$.  Formula (8.23)
then gives
\[
\frac{\partial^2f}{\partial x^i\partial x^j}(0)=h_{ij}(0).
\]
Taylor's theorem yields (8.24).
\end{proof}

If a smooth ambient field $N=N^i\partial_i$ restricts to a unit normal along
$M$, the Weingarten equation gives, for
$X=X^j\partial_j\in TM$,
\begin{equation}
sX=-\overline\nabla_XN
=-\sum_{i,j=1}^{n+1}X^j(\partial_jN^i)\partial_i. \tag{8.25}
\end{equation}
If $M$ is locally a regular level set of $F$, one may take
\[
N=\frac{\operatorname{grad}F}{|\operatorname{grad}F|}.
\]

\begin{example}[Shape Operators of Spheres]
\label{ex:shape-operators-spheres}
\dcref{ch8:8.25}
\uses{prop:existence-local-defining-functions}
For $F(x)=|x|^2$, the outward unit normal on $S^n(R)$ is
\[
N=\frac1R\sum_{i=1}^{n+1}x^i\partial_i.
\]
Therefore
\[
sX=-\frac1R\sum_{i,j=1}^{n+1}X^j(\partial_jx^i)\partial_i
=-\frac1RX.
\]
Thus every principal curvature and the mean curvature equal $-1/R$, while
the Gaussian curvature is $(-1/R)^n$ for this orientation.
\end{example}

For a parametrized surface in $\mathbb R^3$, a unit normal is
\begin{equation}
N=\frac{X_1\times X_2}{|X_1\times X_2|}. \tag{8.26}
\end{equation}

\section{The Gaussian Curvature of a Surface Is Intrinsic}

\begin{exercise}
\label{exer:plane-cylinder-isometric}
\dcref{ch8:8.26}
Let $M_1=\{z=0\}\subset\mathbb R^3$ and
$M_2=\{x^2+y^2=1\}\subset\mathbb R^3$.  Show that they are locally
isometric, although $M_1$ has mean curvature $0$ and $M_2$ has mean curvature
$\pm\tfrac12$, according to the normal orientation.
\end{exercise}

\begin{theorem}[Gauss's Theorema Egregium]
\label{thm:theorema-egregium}
\dcref{ch8:8.27}
\uses{cor:curvature-tensor-dim-2}
For an embedded Riemannian surface $(M,g)\subset\mathbb R^3$, the Gaussian
curvature satisfies
\[
K=\frac12S.
\]
Consequently $K$ is an intrinsic local-isometry invariant.
\end{theorem}

\begin{proof}
\sketch
At $p$, choose an orthonormal basis $(b_1,b_2)$ of $T_pM$.  The shape
operator and $h$ have the same matrix, and the Gauss equation gives
\[
\mathrm{Rm}_p(b_1,b_2,b_2,b_1)
=h_{11}h_{22}-h_{12}h_{21}=\det(h_{ij})=K(p).
\]
In dimension two,
$\mathrm{Rm}=\tfrac14Sg\owedge g$, so evaluation on
$(b_1,b_2,b_2,b_1)$ gives $\mathrm{Rm}_p(b_1,b_2,b_2,b_1)=\tfrac12S(p)$.
\end{proof}

For an abstract Riemannian $2$-manifold, define its \emph{Gaussian curvature}
by $K=\tfrac12S$.  This agrees with the determinant of the shape operator for
embedded surfaces in $\mathbb R^3$.

\begin{corollary}[Curvature Tensors of a Riemannian Surface]
\label{cor:riemann-two-manifold-curvatures}
\dcref{ch8:8.28}
\uses{cor:curvature-tensor-dim-2}
For every Riemannian $2$-manifold,
\[
\mathrm{Rm}=\frac12K g\owedge g,
\qquad
\mathrm{Rc}=Kg,
\qquad
S=2K.
\]
\end{corollary}

\section{Sectional Curvatures}

Let $(M,g)$ be a Riemannian $n$-manifold with $n\geq2$, let $p\in M$, and
let $V\subseteq T_pM$ be a star-shaped neighborhood on which $\exp_p$ is a
diffeomorphism onto $U$.  For a $2$-plane $\Pi\subseteq T_pM$, the
\emph{plane section}
\[
S_\Pi=\exp_p(\Pi\cap V)
\]
is an embedded surface through $p$ with $T_pS_\Pi=\Pi$.  Its intrinsic
Gaussian curvature at $p$ is the \emph{sectional curvature} $\sec(\Pi)$.
For a basis $(v,w)$ of $\Pi$, write $\sec(v,w)=\sec(\Pi)$.

For vectors in an inner-product space, set
\begin{equation}
|v\wedge w|=\sqrt{|v|^2|w|^2-\langle v,w\rangle^2}. \tag{8.27}
\end{equation}
This quantity is positive exactly when $v,w$ are independent and equals $1$
for an orthonormal pair.

\begin{proposition}[Formula for the Sectional Curvature]
\label{prop:sectional-curvature-formula}
\dcref{ch8:8.29}
If $v,w\in T_pM$ are linearly independent, then
\[
\sec(v,w)=\frac{\mathrm{Rm}_p(v,w,w,v)}{|v\wedge w|^2}.
\]
\end{proposition}

\begin{proof}
\sketch
Let $\Pi=\operatorname{span}(v,w)$.  If $z\in\Pi$, the ambient geodesic
$\gamma_z(t)=\exp_p(tz)$ lies in $S_\Pi$ near $0$.  The Gauss formula along
$\gamma_z$ gives
\[
0=D_t\gamma_z'=\overline D_t\gamma_z'+\mathrm{II}(\gamma_z',\gamma_z'),
\]
whose tangent and normal terms vanish separately.  Hence
$\mathrm{II}_p(z,z)=0$ for all $z\in\Pi$, and polarization gives
$\mathrm{II}_p=0$.  The Gauss equation therefore identifies the ambient and
surface curvature tensors on $\Pi$ at $p$.  For an orthonormal basis
$(b_1,b_2)$ of $\Pi$,
\[
\sec(\Pi)=\mathrm{Rm}_p(b_1,b_2,b_2,b_1).
\]

Apply Gram--Schmidt to an arbitrary basis:
\[
b_1=\frac v{|v|},
\qquad
b_2=\frac{w-\langle w,b_1\rangle b_1}
{|w-\langle w,b_1\rangle b_1|}.
\]
The curvature symmetries give
\begin{equation}
\sec(\Pi)=
\frac{\mathrm{Rm}_p(v,w,w,v)}
{|v|^2|w-\langle w,b_1\rangle b_1|^2}. \tag{8.29}
\end{equation}
The denominator is
$|v|^2|w|^2-\langle v,w\rangle^2=|v\wedge w|^2$.
\end{proof}

\begin{exercise}
\label{exer:metric-scaling-sectional-curvature}
\dcref{ch8:8.30}
\uses{thm:conformal-transformation-curvature}
If $\widetilde g=\lambda g$ for a positive constant $\lambda$, use the
conformal curvature transformation to prove
\[
\sec_{\widetilde g}(\Pi)=\lambda^{-1}\sec_g(\Pi)
\]
for every tangent $2$-plane $\Pi$.
\end{exercise}

\begin{proposition}[Sectional Curvature Determines the Curvature Tensor]
\label{prop:sectional-curvature-determines-curvature-tensor}
\dcref{ch8:8.31}
Let $R_1,R_2$ be algebraic curvature tensors on a finite-dimensional
inner-product space $V$.  If
\[
\frac{R_1(v,w,w,v)}{|v\wedge w|^2}
=\frac{R_2(v,w,w,v)}{|v\wedge w|^2}
\]
for every linearly independent $v,w$, then $R_1=R_2$.
\end{proposition}

\begin{proof}
\sketch
Set $D=R_1-R_2$.  Then $D(v,w,w,v)=0$ for all $v,w$, by the hypothesis for
independent pairs and by curvature symmetries for dependent pairs.  Expanding
$D(v+w,x,x,v+w)=0$ gives
\[
D(v,x,x,w)=0.
\]
Expanding $D(v,x+u,x+u,w)=0$ then gives
\[
D(v,x,u,w)+D(v,u,x,w)=0,
\]
so $D$ is antisymmetric in every adjacent pair.  The algebraic Bianchi
identity now yields
\[
0=D(v,w,x,u)+D(w,x,v,u)+D(x,v,w,u)=3D(v,w,x,u),
\]
hence $D=0$.
\end{proof}

\begin{proposition}[Geometric Interpretation of Ricci and Scalar Curvatures]
\label{prop:geometric-interpretation-ricci-scalar-curvatures}
\dcref{ch8:8.32}
\uses{lem:symmetric-bilinear-form-unit-vectors}
Let $(M,g)$ be a Riemannian $n$-manifold and $p\in M$.
\begin{enumerate}
\item[(a)] If $v$ is unit and $(b_1,\ldots,b_n)$ is orthonormal with $b_1=v$,
then
\[
\mathrm{Rc}_p(v,v)=\sum_{k=2}^n\sec(b_1,b_k).
\]
\item[(b)] For any orthonormal basis,
\[
S(p)=\sum_{j\ne k}\sec(b_j,b_k),
\]
the sum over ordered pairs of distinct basis vectors.
\end{enumerate}
\end{proposition}

\begin{proof}
\sketch
Tracing the curvature tensor in the adapted orthonormal basis gives
\[
\mathrm{Rc}_p(v,v)
=\sum_{k=1}^n\mathrm{Rm}_p(b_k,b_1,b_1,b_k)
=\sum_{k=2}^n\sec(b_1,b_k).
\]
Tracing once more gives
\[
S(p)=\sum_{j=1}^n\mathrm{Rc}_p(b_j,b_j)
=\sum_{j\ne k}\sec(b_j,b_k).
\]
\end{proof}

Thus a common strict or weak sign for all sectional curvatures gives the same
sign for Ricci and scalar curvature.  The convention used here is
$\sec(v,w)=\mathrm{Rm}(v,w,w,v)/|v\wedge w|^2$; reversing the curvature-tensor
sign convention requires the corresponding reversed sign/order in this
formula.

\section{Sectional Curvatures of the Model Spaces}

A Riemannian manifold has \emph{constant sectional curvature} if all tangent
$2$-planes at all points have the same sectional curvature.

\begin{lemma}[Frame Homogeneity Implies Constant Sectional Curvature]
\label{lem:frame-homogeneous-constant-sectional-curvature}
\dcref{ch8:8.33}
Every frame-homogeneous Riemannian manifold has constant sectional curvature.
\end{lemma}

\begin{proof}
\sketch
Frame homogeneity supplies an isometry carrying any chosen tangent $2$-plane
to any other; sectional curvature is invariant under isometries.
\end{proof}

\begin{theorem}[Sectional Curvatures of the Model Spaces]
\label{thm:sectional-curvatures-model-spaces}
\dcref{ch8:8.34}
\uses{ex:shape-operators-spheres}
The model spaces have constant sectional curvatures
\begin{enumerate}
\item[(a)] $0$ on $(\mathbb R^n,\bar g)$;
\item[(b)] $1/R^2$ on $(\mathbb S^n(R),\bar g_R)$;
\item[(c)] $-1/R^2$ on $(\mathbb H^n(R),\bar g_R)$.
\end{enumerate}
\end{theorem}

\begin{proof}
\sketch
Euclidean space has zero curvature tensor, hence zero sectional curvature.
For the round sphere, the shape operator is $-(1/R)\operatorname{Id}$, so the
Gauss equation gives sectional curvature $(-1/R)^2=1/R^2$; frame homogeneity
makes this value constant.
\end{proof}

For each $c\in\mathbb R$, exactly one of these model families has constant
sectional curvature $c$, with its radius determined by $|c|^{-1/2}$ when
$c\ne0$.

\begin{exercise}
\label{exer:constant-sectional-curvature-rp}
\dcref{ch8:8.35}
\uses{ex:real-projective-space-metric}
Show that the metric on $\mathbb{RP}^n$ from Example 2.34 has constant positive
sectional curvature.
\end{exercise}

\begin{proposition}[Metrics of Constant Sectional Curvature]
\label{prop:constant-sectional-curvature}
\dcref{ch8:8.36}
A Riemannian metric $g$ has constant sectional curvature $c$ if and only if
\[
\mathrm{Rm}=\frac12c\,g\owedge g.
\]
In this case,
\[
\mathrm{Rc}=(n-1)cg,
\qquad
S=n(n-1)c,
\]
\[
R(v,w)x=c\bigl(\langle w,x\rangle v-\langle v,x\rangle w\bigr),
\]
and, in any basis,
\[
R_{ijkl}=c(g_{il}g_{jk}-g_{ik}g_{jl}),
\qquad
R_{ij}=(n-1)cg_{ij}.
\]
\end{proposition}
\section{Problems}

\begin{exercise}
\label{prob:graph-submanifold-shape}
\dcref{ch8:8-1}
\uses{ex:metrics-in-graph-coordinates}

Let $U\subseteq\mathbb R^n$ be open, let $f:U\to\mathbb R$ be smooth, and
give its graph
$M=\{(x,f(x)):x\in U\}\subseteq\mathbb R^{n+1}$ the induced metric and upward
unit normal.
\begin{enumerate}
\item[(a)] Compute the graph-coordinate components of the shape operator in
terms of the partial derivatives of $f$.
\item[(b)] For $f(x)=|x|^2$, compute all principal curvatures of the resulting
$n$-dimensional paraboloid.
\end{enumerate}
\end{exercise}

\begin{exercise}
\label{prob:hypersurface-defining-function}
\dcref{ch8:8-2}

Let $(M,g)$ be an embedded Riemannian hypersurface of
$(\overline M,\overline g)$, let $F$ be a local defining function for $M$, and
set $N=\operatorname{grad}F/|\operatorname{grad}F|$.
\begin{enumerate}
\item[(a)] Prove that the scalar second fundamental form with respect to $N$
is
\[
h(X,Y)=-\frac{\overline\nabla^2F(X,Y)}{|\operatorname{grad}F|}
\qquad(X,Y\in\mathfrak X(M)).
\]
\item[(b)] If $n=\dim M$, prove
\[
H=-\frac1n\operatorname{div}_{\overline g}
\left(\frac{\operatorname{grad}F}{|\operatorname{grad}F|}\right).
\]
\emph{Hint.} First prove: if $V$ is a finite-dimensional inner product space,
$w\in V$ is unit, and $A:V\to V$ maps $w^\perp$ into itself, then
\[
\operatorname{tr}(A|_{w^\perp})=\operatorname{tr}B,
\qquad
Bx=Ax-\langle x,w\rangle Aw.
\]
\end{enumerate}
\end{exercise}

\begin{exercise}
\label{prob:surface-revolution-shape-principal-curvatures}
\dcref{ch8:8-3}
\uses{ex:surfaces-of-revolution}

Let $C\subseteq\{(r,z):r>0\}$ be an embedded smooth curve, let $S_C$ be its
surface of revolution, let $\gamma(t)=(a(t),b(t))$ be a unit-speed local
parametrization of $C$, and let $X$ be the parametrization in (2.11).
\begin{enumerate}
\item[(a)] Compute the shape operator and principal curvatures in terms of
$a,b$, and prove that the principal directions are tangent to the meridians
and latitude circles.
\item[(b)] Prove that the Gaussian curvature at $X(t,\theta)$ is
$-a''(t)/a(t)$.
\end{enumerate}
\end{exercise}

\begin{exercise}
\label{prob:surface-revolution-constant-gaussian-curvature}
\dcref{ch8:8-4}
\uses{cor:local-uniqueness-constant-curvature}

Construct a surface of revolution in $\mathbb R^3$ with constant Gaussian
curvature $1$ but nonconstant principal curvatures. Together with
Corollary~10.15, this yields locally isometric nonflat Euclidean hypersurfaces
with different principal curvatures.
\end{exercise}

\begin{exercise}
\label{prob:paraboloid-isotropic-one-point}
\dcref{ch8:8-5}

For the paraboloid
$S=\{(x,y,z)\in\mathbb R^3:z=x^2+y^2\}$ with its induced metric, prove that
$S$ is isotropic at exactly one point.
\end{exercise}

\begin{exercise}
\label{prob:geodesic-curvature-non-unit-speed-curve}
\dcref{ch8:8-6}

Let $\gamma:I\to(M,g)$ be a regular curve in a Riemannian manifold, with no
unit-speed assumption. Prove
\[
\kappa(t)=
\frac{|\gamma'(t)\times D_t\gamma'(t)|}{|\gamma'(t)|^3},
\]
where the numerator uses the norm in (8.27). In Euclidean $\mathbb R^3$,
deduce
\[
\kappa(t)=
\frac{|\gamma'(t)\times\gamma''(t)|}{|\gamma'(t)|^3}.
\]
\end{exercise}

\begin{exercise}
\label{prob:catenoid-minimal-surface}
\dcref{ch8:8-7}

For $w>0$, revolve
$\gamma(t)=(w\cosh(t/w),t)$ about the $z$-axis to obtain the
\emph{catenoid} $M_w\subseteq\mathbb R^3$. Prove that every $M_w$ is minimal.
\end{exercise}

\begin{exercise}
\label{prob:catenoid-boundary-minimal}
\dcref{ch8:8-8}
\uses{prob:catenoid-minimal-surface}

For $h>0$, let
\[
C_h=\{(x,y,z)\in\mathbb R^3:x^2+y^2=1,\ z=\pm h\}.
\]
Let $c>0$ be the unique solution of $c\tanh c=1$ and set
$h_0=1/\sinh c$. Prove that:
\begin{enumerate}
\item if $h<h_0$, exactly two values $0<w_1<w_2$ make the portions of
$M_{w_i}$ in $|z|\leq h$ minimal surfaces with boundary $C_h$;
\item if $h=h_0$, exactly one positive $w$ has this property;
\item if $h>h_0$, no positive $w$ has this property.
\end{enumerate}
\end{exercise}

\begin{exercise}
\label{prob:gauss-map}
\dcref{ch8:8-9}

Let $M\subseteq\mathbb R^{n+1}$ be a Riemannian hypersurface with smooth unit
normal $N$. Its Gauss map is
\[
\nu:M\to S^n,\qquad \nu(p)=N_p.
\]
If $\mathring g$ is the round metric on $S^n$, prove
\[
\nu^*dV_{\mathring g}=(-1)^nK\,dV_g,
\]
where $K$ is the Gaussian curvature of $M$.
\end{exercise}

\begin{exercise}
\label{prob:product-metric-submanifolds}
\dcref{ch8:8-10}
\uses{exer:prove-lemma-2-11,prob:product-metric-curvature}

Let $g=g_1\oplus g_2$ be the product metric on $M_1\times M_2$ from (2.12).
\begin{enumerate}
\item[(a)] For $p_i\in M_i$, prove that
$M_1\times\{p_2\}$ and $\{p_1\}\times M_2$ are totally geodesic.
\item[(b)] Prove $\sec(\Pi)=0$ whenever
$\Pi\subseteq T_{(p_1,p_2)}(M_1\times M_2)$ is spanned by
$v_1\in T_{p_1}M_1$ and $v_2\in T_{p_2}M_2$.
\item[(c)] Prove that the product has nonnegative sectional curvature when
both factors do, and nonpositive sectional curvature when both factors do.
\item[(d)] Suppose $(M_i,g_i)$ has constant sectional curvature $c_i$ and
$n_i=\dim M_i$. Prove that the product is Einstein exactly when
\[
c_1(n_1-1)=c_2(n_2-1),
\]
and has constant sectional curvature exactly when
$n_1=n_2=1$ or $c_1=c_2=0$.
\end{enumerate}
\end{exercise}

\begin{exercise}
\label{prob:surface-affine-plane-curvature}
\dcref{ch8:8-11}

Let $M\subseteq\mathbb R^3$ be a smooth surface, let $p\in M$, and let
$N_p$ be normal to $M$ at $p$. If an affine plane
$\Pi\subseteq\mathbb R^3$ contains $p$ and is parallel to $N_p$, prove that
some neighborhood $U$ of $p$ satisfies:
$M\cap\Pi\cap U$ is a smooth embedded curve in $\Pi$, and the principal
curvatures of $M$ at $p$ are the minimum and maximum signed Euclidean
curvatures of these curves as $\Pi$ varies.
\end{exercise}

\begin{exercise}
\label{prob:riemannian-immersion-bounded-sectional-curvatures}
\dcref{ch8:8-12}

Let $\pi:(\widetilde M,\widetilde g)\to(M,g)$ be a Riemannian submersion.
If every sectional curvature of $(\widetilde M,\widetilde g)$ is at least
$c$, use O'Neill's formula from Problem~7-14 to prove that every sectional
curvature of $(M,g)$ is at least $c$.
\end{exercise}

\begin{exercise}
\label{prob:complex-projective-space-curvature}
\dcref{ch8:8-13}

Let $p:S^{2n+1}\to\mathbb{CP}^n$ be the Riemannian submersion of
Example~2.30. Identify $\mathbb C^{n+1}$ with $\mathbb R^{2n+2}$ using
$z^j=x^j+iy^j$.
\begin{enumerate}
\item[(a)] Prove that
\[
S=x^j\frac{\partial}{\partial y^j}
-y^j\frac{\partial}{\partial x^j}
\]
is tangent to $S^{2n+1}$ and spans the vertical space $V_z$ at each
$z\in S^{2n+1}$, with $j$ summed from $1$ to $n+1$.
\item[(b)] For horizontal vector fields $W,Z$ on $S^{2n+1}$, prove
\[
[W,Z]^V=-d\omega(W,Z)S=2\langle W,JZ\rangle S,
\]
where
\[
\omega=S^\flat=\sum_j(x^j\,dy^j-y^j\,dx^j)
\]
and the real-linear orthogonal map $J$ is
\[
J\left(a^j\frac{\partial}{\partial x^j}
+b^j\frac{\partial}{\partial y^j}\right)
=a^j\frac{\partial}{\partial y^j}
-b^j\frac{\partial}{\partial x^j}.
\]
Thus $J$ is complex multiplication by $i$ and $J^2=-\operatorname{Id}$.
\item[(c)] Use O'Neill's formula from Problem~7-14 to prove
\begin{align*}
Rm(w,x,y,z)
&=\langle\widetilde w,\widetilde z\rangle
  \langle\widetilde x,\widetilde y\rangle
-\langle\widetilde w,\widetilde y\rangle
  \langle\widetilde x,\widetilde z\rangle\\
&\quad-2\langle\widetilde w,J\widetilde x\rangle
  \langle\widetilde y,J\widetilde z\rangle
-\langle\widetilde w,J\widetilde y\rangle
  \langle\widetilde x,J\widetilde z\rangle\\
&\quad+\langle\widetilde w,J\widetilde z\rangle
  \langle\widetilde x,J\widetilde y\rangle
\end{align*}
for $q\in\mathbb{CP}^n$ and $w,x,y,z\in T_q\mathbb{CP}^n$, where the tilded
vectors are horizontal lifts to any
$\widetilde q\in p^{-1}(q)\subseteq S^{2n+1}$.
\item[(d)] For orthonormal $w,x\in T_q\mathbb{CP}^n$, prove
\[
\sec(w,x)=1+3\langle\widetilde w,J\widetilde x\rangle^2.
\]
\end{enumerate}
\end{exercise}

\begin{exercise}
\label{prob:einstein-sectional-curvature}
\dcref{ch8:8-14}

Prove that a Riemannian $3$-manifold is Einstein if and only if it has
constant sectional curvature.
\end{exercise}

\begin{exercise}
\label{prob:homogeneous-isotropic-3-manifold}
\dcref{ch8:8-15}
\uses{prob:complex-projective-space-curvature}

Prove that every homogeneous and isotropic Riemannian $3$-manifold has
constant sectional curvature, and give a counterexample in dimension $4$.
\emph{Hint.} Use Problem~8-13.
\end{exercise}

\begin{exercise}
\label{prob:berger-metric-su2}
\dcref{ch8:8-16}

For $a>0$, let $g_a$ be the Berger metric on $\operatorname{SU}(2)$ from
Problem~3-10. Compute the sectional curvatures of the planes spanned by
$(X,Y)$, $(Y,Z)$, and $(Z,X)$. Prove that when $a\neq1$,
$(\operatorname{SU}(2),g_a)$ is homogeneous but nowhere isotropic.
\end{exercise}

\begin{exercise}
\label{prob:lie-group-bi-invariant-metric-ch8}
\dcref{ch8:8-17}

Let $G$ be a Lie group with a bi-invariant metric $g$.
\begin{enumerate}
\item[(a)] For orthonormal $X,Y\in\operatorname{Lie}(G)$, prove
\[
\sec(X_p,Y_p)=\frac14|[X,Y]|^2
\qquad(p\in G),
\]
and deduce that all sectional curvatures are nonnegative.
\item[(b)] Prove that every Lie subgroup of $G$ is totally geodesic.
\item[(c)] If $G$ is connected, prove that $G$ is flat exactly when it is
abelian.
\end{enumerate}
\end{exercise}

\begin{exercise}
\label{prob:convex-hypersurface}
\dcref{ch8:8-18}

Let $(M,g)$ be a Riemannian hypersurface in
$(\overline M,\overline g)$ with unit normal $N$. Assume $M$ is
\emph{convex with respect to $N$}, meaning
$h(v,v)\leq0$ for every $v\in TM$. If the ambient sectional curvatures are
bounded below by $c$, prove that the sectional curvatures of $M$ are also
bounded below by $c$.
\end{exercise}

\begin{exercise}
\label{prob:lorentz-hypersurface-formulas}
\dcref{ch8:8-19}
\uses{thm:fundamental-equations-hypersurface}

Let $(M,g)$ be a Riemannian hypersurface in an $(n+1)$-dimensional Lorentz
manifold $(\overline M,\overline g)$ with smooth unit normal $N$. Define
$h$ and $s$ by
\[
\mathrm{II}(X,Y)=h(X,Y)N,
\qquad
\langle sX,Y\rangle=h(X,Y).
\]
Prove the following Lorentzian analogues of Theorem~8.13:
\begin{enumerate}
\item[(a)] \emph{Gauss formula:}
$\overline\nabla_XY=\nabla_XY+h(X,Y)N$.
\item[(b)] \emph{Gauss formula along a curve:}
$\overline D_tX=D_tX+h(\gamma',X)N$.
\item[(c)] \emph{Weingarten equation:}
$\overline\nabla_XN=sX$.
\item[(d)] \emph{Gauss equation:}
\[
\overline{Rm}(W,X,Y,Z)
=\left(Rm+\frac12h\owedge h\right)(W,X,Y,Z).
\]
\item[(e)] \emph{Codazzi equation:}
$\overline{Rm}(W,X,Y,N)=-(Dh)(Y,W,X)$.
\end{enumerate}
\end{exercise}

\begin{exercise}
\label{prob:lorentz-einstein-field-equation}
\dcref{ch8:8-20}
\uses{prob:lorentz-hypersurface-formulas}

Let $(\overline M,\overline g)$ be an $(n+1)$-dimensional Lorentz manifold
satisfying
\[
\overline{\mathrm{Rc}}-\frac12\overline S\,\overline g
+\Lambda\overline g=T,
\]
where $\Lambda$ is constant and $T$ is a smooth symmetric $2$-tensor. Let
$(M,g)$ be a Riemannian hypersurface with scalar second fundamental form $h$
as in Problem~8-19. Prove the Einstein constraint equations
\[
S-2\Lambda-|h|_g^2+(\operatorname{tr}_gh)^2=2\rho,
\]
\[
\operatorname{div}h-d(\operatorname{tr}_gh)=J,
\]
where $\rho=T(N,N)$ and $J(X)=T(N,X)$. Here $S$ is the scalar curvature of
$g$, and $\operatorname{div}h$ is the divergence from Problem~5-16. When
$J=\rho=0$, these are the \emph{vacuum Einstein constraint equations}.
\end{exercise}

\begin{exercise}[The associated quadratic form]
\label{prob:associated-quadratic-form}
\dcref{ch8:8-21}

For a linear endomorphism $A:\mathbb R^n\to\mathbb R^n$, define its
\emph{associated quadratic form} by $Q(x)=\langle Ax,x\rangle$. Prove
\[
\int_{S^{n-1}}Q\,dV_g
=\frac1n(\operatorname{tr}A)\operatorname{Vol}(S^{n-1}).
\]
\emph{Hint.} For $i\neq j$, prove
$\int_{S^{n-1}}x^ix^j\,dV_g=0$ using the isometry
\[
(x^1,\ldots,x^i,\ldots,x^j,\ldots,x^n)
\longmapsto
(x^1,\ldots,x^j,\ldots,-x^i,\ldots,x^n).
\]
Then compute $\int_{S^{n-1}}(x^i)^2\,dV_g$ from
$\sum_i(x^i)^2=1$.
\end{exercise}

\begin{exercise}[Ricci curvature without choice of basis]
\label{prob:ricci-curvature-without-basis}
\dcref{ch8:8-22}
\uses{prob:associated-quadratic-form}

Let $(M,g)$ be a Riemannian $n$-manifold, $p\in M$, and
$v\in T_pM$ a unit vector. Prove
\[
\operatorname{Rc}_p(v,v)
=\frac{n-1}{\operatorname{Vol}(S^{n-2})}
\int_{w\in S_v^\perp}\sec(v,w)\,dV_{\widehat g},
\]
where $S_v^\perp$ is the sphere of unit vectors in $T_pM$ orthogonal to $v$,
and $\widehat g$ is induced on it by the flat metric $g_p$.
\emph{Hint.} Use Problem~8-21.
\end{exercise}

\begin{exercise}
\label{prob:lie-group-isometric-action-dimension}
\dcref{ch8:8-23}
\uses{prob:lie-group-dimension-bound}

Let a Lie group $G$ act effectively and isometrically on a connected
Riemannian $n$-manifold $(M,g)$. Given the bound
$\dim G\leq n(n+1)/2$ from Problem~5-12, prove that equality implies constant
sectional curvature.
\end{exercise}

\begin{exercise}
\label{prob:k-point-homogeneous-ch8}
\dcref{ch8:8-24}
\uses{prob:k-point-homogeneous}

For a connected Riemannian manifold $(M,g)$, prove:
\begin{enumerate}
\item[(a)] homogeneity implies constant scalar curvature;
\item[(b)] $2$-point homogeneity implies that $g$ is Einstein;
\item[(c)] $3$-point homogeneity implies constant sectional curvature.
\end{enumerate}
\end{exercise}

\begin{exercise}
\label{prob:weyl-tensor-product-metric}
\dcref{ch8:8-25}
\uses{prob:product-metric-curvature}

For $i=1,2$, let $(M_i,g_i)$ have dimension $n_i\geq2$ and constant
sectional curvature $c_i$. On $M=M_1\times M_2$, set
$g=g_1\oplus g_2$, $n=n_1+n_2$, and $h_i=\pi_i^*g_i$. Prove
\[
W=\frac{c_1+c_2}{2(n-1)(n-2)}
\left(
n_2(n_2-1)h_1\owedge h_1
-2(n_1-1)(n_2-1)h_1\owedge h_2
+n_1(n_1-1)h_2\owedge h_2
\right).
\]
Deduce that $(M,g)$ is locally conformally flat if and only if
$c_2=-c_1$.
\end{exercise}

\begin{exercise}
\label{prob:weyl-tensor-4d}
\dcref{ch8:8-26}

Let $(M,g)$ be a Riemannian $4$-manifold, let $p\in M$, and let $(b_i)$ be
an orthonormal basis of $T_pM$. Prove
\[
W_{1221}
=\frac13(k_{12}+k_{34})
-\frac16(k_{13}+k_{14}+k_{23}+k_{24}),
\]
where $k_{ij}$ is the sectional curvature of
$\operatorname{span}(b_i,b_j)$.
\end{exercise}

\begin{exercise}
\label{prob:fubini-study-not-conformal-flat}
\dcref{ch8:8-27}
\uses{prob:complex-projective-space-curvature,prob:weyl-tensor-4d}

Prove that the Fubini--Study metric on $\mathbb{CP}^2$ is not locally
conformally flat. \emph{Hint.} Use Problems~8-13 and~8-26.
\end{exercise}

\begin{exercise}
\label{prob:hyperbolic-constant-curvature}
\dcref{ch8:8-28}
\uses{thm:sectional-curvatures-model-spaces,thm:gauss-equation,prob:lorentz-hypersurface-formulas}

Complete the proof of Theorem~8.34 by proving in two ways that
$\mathbb H^n(R)$ has constant sectional curvature $-1/R^2$.
\begin{enumerate}
\item[(a)] In the hyperboloid model, compute the second fundamental form of
$\mathbb H^n(R)\subseteq\mathbb R^{n,1}$ at $(0,\ldots,0,R)$ and apply either
the Gauss equation or Problem~8-19.
\item[(b)] In the Poincar\'e ball model, use the conformal curvature
transformation (7.44) to compute the Riemann tensor at the origin.
\end{enumerate}
\end{exercise}

\begin{exercise}
\label{prob:curvature-tensors-constant-curvature}
\dcref{ch8:8-29}
\uses{prop:constant-sectional-curvature}

Prove Proposition~8.36 on curvature tensors of constant-curvature spaces.
\end{exercise}

\begin{exercise}
\label{prob:ricci-tensor-hypersurface}
\dcref{ch8:8-30}

Let $M\subseteq\mathbb R^{n+1}$ be a hypersurface with induced metric. Prove
\[
\mathrm{Rc}(v,w)=\langle nHsv-s^2v,w\rangle,
\]
where $H$ and $s$ are its mean curvature and shape operator, and
$s^2v=s(sv)$.
\end{exercise}

\begin{exercise}
\label{prob:k-points-hyperbolic-geodesic}
\dcref{ch8:8-31}

For $1<k\leq n$, prove that any $k$ points of $\mathbb H^n(R)$ lie in a
totally geodesic $(k-1)$-dimensional submanifold isometric to
$\mathbb H^{k-1}(R)$.
\end{exercise}

\begin{exercise}
\label{prob:lie-group-fixed-point-geodesic}
\dcref{ch8:8-32}

Let a Lie group $G$ act smoothly and isometrically on a Riemannian manifold
$(M,g)$, and let
\[
S=\{p\in M:\varphi(p)=p\text{ for every }\varphi\in G\}.
\]
Prove that every connected component of $S$ is a smoothly embedded totally
geodesic submanifold. Different components are allowed to have different
dimensions.
\end{exercise}

\begin{exercise}
\label{prob:curvature-operator}
\dcref{ch8:8-33}
\uses{prob:complex-projective-space-curvature}

Let $(M,g)$ be Riemannian, and let $\Lambda^2(TM)$ be its bundle of
alternating contravariant $2$-tensors.
\begin{enumerate}
\item[(a)] Prove that there is a unique fiber metric on $\Lambda^2(TM)$ whose
norm satisfies
\[
|w\wedge x|^2=|w|^2|x|^2-\langle w,x\rangle^2.
\]
\item[(b)] Prove that there is a unique bundle endomorphism
$\mathcal R:\Lambda^2(TM)\to\Lambda^2(TM)$ satisfying
\begin{equation}
\langle\mathcal R(w\wedge x),y\wedge z\rangle=-Rm(w,x,y,z)
\tag{8.30}
\end{equation}
for tangent vectors at a common point. This is the \emph{curvature operator}.
\item[(c)] Define positive and negative curvature operator by positive and
negative definiteness of $\mathcal R$. Prove that they imply positive and
negative sectional curvature, respectively.

\emph{Sign convention.} The minus sign in (8.30) is required by the curvature
convention used here; it is omitted when the Riemann tensor is defined with
the opposite sign.
\item[(d)] Show that the converse need not hold. On
$\mathbb{CP}^2$ with the Fubini--Study metric, compute
\[
\left\langle
\mathcal R(w\wedge x+y\wedge z),w\wedge x+y\wedge z
\right\rangle,
\]
where $w,x,y,z$ are orthonormal and their horizontal lifts satisfy
$\widetilde y=J\widetilde w$ and $\widetilde z=J\widetilde x$.
\end{enumerate}
\end{exercise}
